\chapter{Probabilistic Data Structures}\label{ch19:probabilistic-data-structures}

\section{Introduction}

Not all problems require exact answers. When processing massive data streams, probabilistic data structures\index{probabilistic data structure} provide memory-efficient approximate answers with bounded error.

> \textbf{Complete compilable implementations of the data structures in this chapter are in \texttt{code/}.}

\section{Bloom Filters}

A \textbf{Bloom filter}\index{Bloom filter} tests set membership with no false negatives and a controllable false positive rate. It uses k hash functions and a bit array of size m.

\subsection{Operations}

\begin{lstlisting}[language=C++]
class bloom_filter {
public:
    bloom_filter(size_t expected_elements, double false_positive_rate = 0.01) {
        // Optimal size and hash count
        m_ = -expected_elements * std::log(false_positive_rate) 
             / (std::log(2) * std::log(2));
        k_ = std::max(1, static_cast<int>(std::round(
             (m_ / expected_elements) * std::log(2))));
        bits_.resize(m_, false);
    }

    void insert(std::string_view key) {
        auto [h1, h2] = hash(key);
        for (int i = 0; i < k_; ++i) {
            size_t pos = (h1 + i * h2) % m_;
            bits_[pos] = true;
        }
    }

    bool contains(std::string_view key) const {
        auto [h1, h2] = hash(key);
        for (int i = 0; i < k_; ++i) {
            size_t pos = (h1 + i * h2) % m_;
            if (!bits_[pos]) return false;  // definitely not present
        }
        return true;  // probably present
    }

private:
    // Double hashing
    std::pair<size_t, size_t> hash(std::string_view key) const {
        std::hash<std::string_view> h;
        size_t hval = h(key);
        return {hval, hval >> 16 | (hval << 16)};
    }

    size_t m_;  // bits in filter
    int k_;     // number of hash functions
    std::vector<bool> bits_;
};
\end{lstlisting}

\textbf{Properties:}
\begin{itemize}
\item \textbf{No false negatives:} if contains() returns false, the key was definitely not inserted
\item \textbf{False positives:} contains() may return true for a key that was not inserted
\item \textbf{Cannot delete} (without extensions like counting Bloom filters\index{counting Bloom filter})
\end{itemize}

\textbf{Optimal parameters:} m = -n ln P / (ln 2)$^{2}$, k = (m/n) ln 2, where n = expected elements, P = target false positive rate.

\textbf{Application:} Cache filtering, spell checkers, database join optimization, cryptocurrency SPV nodes.

\section{Cuckoo Filters}

A \textbf{cuckoo filter}\index{cuckoo filter} (Fan et al., 2014) is a Bloom filter alternative that supports \textbf{deletion} with better space efficiency for low false positive rates. The name comes from ``cuckoo hashing'' — items that collide are displaced to alternate locations, like a cuckoo bird evicting eggs.

\subsection{Core Idea}

Like a Bloom filter, a cuckoo filter stores fingerprints (short hash values) in a fixed-size array. Each element maps to \textbf{two candidate buckets} (not $k$ buckets like a Bloom filter). When inserting:

\begin{enumerate}
\item Compute the fingerprint $f$ of the key and the two candidate bucket indices $i_1, i_2$
\item If either bucket has an empty slot, place $f$ there
\item If both buckets are full, \textbf{evict} an existing fingerprint to its alternate bucket (the other candidate for that fingerprint)
\item After at most a constant number of evictions (typically 500), declare the filter full and trigger a rebuild
\end{enumerate}

\begin{lstlisting}[language=Java]
  Bucket 7:  [f1] [f2] [  ] [  ]       Bucket 12: [f3] [f4] [f5] [f6]
              |                              |
           i1 = 7                         i2 = 12

  Insert f7 (maps to buckets 7 and 12):
    Both full  ->  evict f3 from bucket 12 to its alternate (bucket 7 or another)

\end{lstlisting}

\subsection{Cuckoo Filter Implementation}

\begin{lstlisting}[language=C++]
class cuckoo_filter {
public:
    cuckoo_filter(size_t capacity, double target_fpr = 0.01)
        : buckets_(capacity) {
        // fingerprint size in bits: ceil(log2(1/fpr)) + 3 for amortized analysis
        fp_bits_ = static_cast<int>(std::ceil(std::log2(1.0 / target_fpr))) + 3;
        fp_mask_ = (1ULL << fp_bits_) - 1;
    }

    void insert(std::string_view key) {
        auto f = fingerprint(key);
        auto [i1, i2] = bucket_indices(key, f);

        // Try to insert into either bucket
        if (try_insert(i1, f)) return;
        if (try_insert(i2, f)) return;

        // Both full -- evict from the first bucket
        for (int n = 0; n < max_kicks; ++n) {
            size_t victim_idx = (n % 2 == 0) ? i1 : i2;
            auto evicted = kickout(victim_idx, f);
            f = evicted;
            // Compute alternate bucket for the evicted fingerprint
            victim_idx = alt_index(victim_idx, evicted);
            if (try_insert(victim_idx, f)) return;
        }
        // If we get here, filter is too full -- rebuild with larger array
    }

    bool contains(std::string_view key) const {
        auto f = fingerprint(key);
        auto [i1, i2] = bucket_indices(key, f);
        return bucket_contains(i1, f) || bucket_contains(i2, f);
    }

    void erase(std::string_view key) {
        auto f = fingerprint(key);
        auto [i1, i2] = bucket_indices(key, f);
        if (remove_from_bucket(i1, f)) return;
        remove_from_bucket(i2, f);
    }

private:
    static constexpr int max_kicks = 500;
    static constexpr int bucket_size_ = 4;

    struct bucket {
        uint64_t fingerprints[4] = {0, 0, 0, 0};
        uint8_t count = 0;
    };

    std::vector<bucket> buckets_;
    int fp_bits_;
    uint64_t fp_mask_;

    uint64_t fingerprint(std::string_view key) const {
        uint64_t h = std::hash<std::string_view>{}(key);
        return (h ^ (h >> 32)) & fp_mask_;
    }

    std::pair<size_t, size_t> bucket_indices(std::string_view key, uint64_t f) const {
        uint64_t h = std::hash<std::string_view>{}(key);
        size_t i1 = h % buckets_.size();
        size_t i2 = alt_index(i1, f);
        return {i1, i2};
    }

    // Partial-key hashing: derive the alternate index from the bucket index
    // and fingerprint alone (no need to re-hash the original key)
    size_t alt_index(size_t index, uint64_t f) const {
        return index ^ static_cast<size_t>(
            ((f >> 16) ^ (f * 0x517cc1b727220a95ULL)) % buckets_.size());
    }

    bool try_insert(size_t idx, uint64_t f) {
        if (buckets_[idx].count < bucket_size_) {
            buckets_[idx].fingerprints[buckets_[idx].count++] = f;
            return true;
        }
        return false;
    }

    uint64_t kickout(size_t idx, uint64_t new_f) {
        int slot = buckets_[idx].count - 1;
        uint64_t old_f = buckets_[idx].fingerprints[slot];
        buckets_[idx].fingerprints[slot] = new_f;
        return old_f;
    }

    bool bucket_contains(size_t idx, uint64_t f) const {
        for (int i = 0; i < buckets_[idx].count; ++i)
            if (buckets_[idx].fingerprints[i] == f) return true;
        return false;
    }

    bool remove_from_bucket(size_t idx, uint64_t f) {
        for (int i = 0; i < buckets_[idx].count; ++i) {
            if (buckets_[idx].fingerprints[i] == f) {
                buckets_[idx].fingerprints[i] =
                    buckets_[idx].fingerprints[--buckets_[idx].count];
                return true;
            }
        }
        return false;
    }
};
\end{lstlisting}

The key insight is \textbf{partial-key hashing}: the alternate bucket index is derived from the current index and the fingerprint alone, without re-hashing the original key. This means evictions don't require knowing the original key — they just swap fingerprints between two buckets.

\subsection{Cuckoo Filter vs Bloom Filter}

\begin{center}
\begin{tabular}{ccccc}
\toprule
Property & Bloom Filter & Cuckoo Filter \\
Space at 1\% FPR & 9.6 bits/elem & 8.0 bits/elem \\
Space at 0.1\% FPR & 14.4 bits/elem & 12.0 bits/elem \\
Deletion & No (without counting) & Yes \\
Lookup speed & $k$ hash probes & 2 memory accesses \\
Cache friendliness & Moderate & High (2 buckets) \\

\bottomrule
\end{tabular}
\end{center}

Cuckoo filters use less space because each bucket stores 4 fingerprints (32+ bits) rather than individual bits — better utilization of cache lines. The trade-off: cuckoo filters can fail to insert when the array is too full, requiring a rebuild.

\textbf{When to use which:}
\begin{itemize}
\item \textbf{Bloom filter:} When deletions are not needed, or when the element count is unknown and may grow without bound
\item \textbf{Cuckoo filter:} When deletions are needed, or when space is tight and the maximum element count is known
\end{itemize}

\section{Count-Min Sketch}

Count approximate frequencies of items in a stream using a 2D array of counters and hash functions.

\begin{lstlisting}[language=C++]
class count_min_sketch {
public:
    count_min_sketch(double epsilon, double delta) {
        // width = ceil(e / epsilon), depth = ceil(ln(1 / delta))
        width_ = static_cast<size_t>(std::ceil(std::exp(1) / epsilon));
        depth_ = static_cast<size_t>(std::ceil(std::log(1.0 / delta)));
        counts_.resize(depth_, std::vector<size_t>(width_, 0));
    }

    void add(std::string_view key, size_t count = 1) {
        for (size_t i = 0; i < depth_; ++i) {
            size_t pos = hash_i(key, i) % width_;
            counts_[i][pos] += count;
        }
    }

    size_t estimate(std::string_view key) const {
        size_t min_val = SIZE_MAX;
        for (size_t i = 0; i < depth_; ++i) {
            size_t pos = hash_i(key, i) % width_;
            min_val = std::min(min_val, counts_[i][pos]);
        }
        return min_val;
    }

private:
    size_t hash_i(std::string_view key, size_t i) const {
        return std::hash<std::string_view>{}(key) + i * 0x9e3779b97f4a7c15ULL;
    }

    size_t width_, depth_;
    std::vector<std::vector<size_t>> counts_;
};
\end{lstlisting}

\section{Heavy Hitters (Misra-Gries)}

Find the most frequent items in a stream — items that appear more than $\frac{n}{k}$ times. The Misra-Gries algorithm uses O(k) counters to identify all heavy hitters.

\subsection{The Idea}

Maintain at most k-1 (item, count) pairs. For each new item:
\begin{itemize}
\item If it matches an existing entry, increment its count.
\item If not and we have fewer than k-1 entries, add it with count 1.
\item If we have k-1 entries, decrement all counts by 1 and remove any entries with count 0.
\end{itemize}

After processing the stream, any item with count $> 0$ is a candidate. Verify candidates by a second pass.

\subsection{Implementation}

\begin{lstlisting}[language=C++]
template <typename Key>
class misra_gries {
public:
    explicit misra_gries(int k) : k_(k) {}

    void add(const Key& key) {
        // If key already tracked, increment
        auto it = counters_.find(key);
        if (it != counters_.end()) {
            it->second += 1;
            return;
        }

        // If room, add new entry
        if (static_cast<int>(counters_.size()) < k_ - 1) {
            counters_[key] = 1;
            return;
        }

        // Decrement all, remove zeros
        for (auto it = counters_.begin(); it != counters_.end(); ) {
            it->second -= 1;
            if (it->second == 0) it = counters_.erase(it);
            else ++it;
        }
    }

    std::vector<std::pair<Key, size_t>> heavy_hitters(size_t min_count) const {
        std::vector<std::pair<Key, size_t>> result;
        for (const auto& [key, count] : counters_) {
            if (count >= min_count) result.push_back({key, count});
        }
        return result;
    }

    const std::unordered_map<Key, size_t>& counters() const { return counters_; }

private:
    int k_;
    std::unordered_map<Key, size_t> counters_;
};
\end{lstlisting}

\subsection{Complexity}

\begin{itemize}
\item \textbf{Space:} O(k) counters
\item \textbf{Update:} O(1) amortized
\item \textbf{Guarantee:} any item with frequency $> \frac{n}{k}$ is among the candidates. False positives possible — verify with a second pass.
\end{itemize}

\subsection{Applications}

\begin{itemize}
\item \textbf{Network security:} detect DDoS attacks by finding IPs sending $> \frac{n}{k}$ requests
\item \textbf{Query analysis:} find most popular search queries at Google/Bing scale
\item \textbf{Log analytics:} top error codes in server logs
\item \textbf{Clickstream analysis:} most visited pages in the last hour
\end{itemize}

\textbf{Relationship to Count-Min Sketch:} Count-Min Sketch\index{Count-Min sketch} gives approximate frequency for any item. Misra-Gries gives exact identification of heavy hitters with less space. For heavy hitters specifically, Misra-Gries is preferred. For arbitrary frequency queries, Count-Min is better.

\section{HyperLogLog}

Estimate the number of distinct elements in a large stream using O(log log n) space.

\begin{lstlisting}[language=C++]
class hyperloglog {
public:
    hyperloglog(int precision = 14) : precision_(precision) {
        m_ = 1 << precision;  // number of registers
        registers_.resize(m_, 0);
    }

    void add(std::string_view key) {
        auto hash = std::hash<std::string_view>{}(key);
        size_t idx = hash >> (64 - precision_);  // first p bits
        uint8_t zeros = count_leading_zeros(hash << precision_);
        registers_[idx] = std::max(registers_[idx], zeros + 1);
    }

    double estimate() const {
        double sum = 0.0;
        size_t zeros = 0;
        for (uint8_t r : registers_) {
            sum += 1.0 / (1 << r);
            if (r == 0) ++zeros;
        }
        double e = alpha_m_ * m_ * m_ / sum;

        // Small range correction
        if (e <= 2.5 * m_) {
            if (zeros > 0) e = m_ * std::log(static_cast<double>(m_) / zeros);
        }
        return e;
    }

private:
    static uint8_t count_leading_zeros(size_t x) {
        return __builtin_clzll(x);
    }

    int precision_;
    size_t m_;
    std::vector<uint8_t> registers_;
    static constexpr double alpha_m_ = 0.7213 / (1 + 1.079 / (1 << 14));
};
\end{lstlisting}

\section{Quantile Estimation}

Estimate the median or any quantile from a data stream without storing all elements. A \textbf{quantile} $q$ is the value below which $q$ fraction of elements fall (e.g., median = 0.5-quantile, P99 = 0.99-quantile).

\subsection{Greenwald-Khanna Sketch}

The GK sketch maintains a set of summary tuples $(v_i, g_i, \Delta_i)$ where $v_i$ is a value, $g_i$ is the maximum gap to the next summary value, and $\Delta_i$ is the maximum gap to the previous summary value.

\begin{lstlisting}[language=C++]
class quantile_estimator {
public:
    explicit quantile_estimator(double epsilon) : epsilon_(epsilon) {}

    void add(double value) {
        ++total_count_;
        summary_.push_back({value, 1, 0});

        // Merge adjacent tuples to maintain size bound
        if (summary_.size() > static_cast<size_t>(1.0 / epsilon_)) {
            merge_tuples();
        }
    }

    double query(double phi) const {
        // Query the phi-quantile (e.g., 0.5 for median)
        if (summary_.empty()) return 0.0;

        double rank = phi * total_count_;
        size_t cumulative = 0;

        for (size_t i = 0; i < summary_.size(); ++i) {
            cumulative += summary_[i].g;
            if (cumulative + summary_[i].delta >= static_cast<size_t>(rank) &&
                cumulative <= static_cast<size_t>(rank)) {
                return summary_[i].value;
            }
        }
        return summary_.back().value;
    }

private:
    struct tuple {
        double value;
        size_t g;      // gap to next
        size_t delta;   // max gap to previous
    };

    void merge_tuples() {
        std::vector<tuple> merged;
        for (size_t i = 0; i + 1 < summary_.size(); ++i) {
            if (summary_[i].g + summary_[i+1].g +
                summary_[i+1].delta <= 2 * epsilon_ * total_count_) {
                // Merge: skip this tuple
                merged.push_back({summary_[i].value,
                    summary_[i].g + summary_[i+1].g,
                    summary_[i+1].delta});
                ++i;  // skip next
            } else {
                merged.push_back(summary_[i]);
            }
        }
        if (!merged.empty()) merged.back().delta = 0;
        summary_ = std::move(merged);
    }

    double epsilon_;
    size_t total_count_ = 0;
    std::vector<tuple> summary_;
};
\end{lstlisting}

\subsection{Complexity}

\begin{itemize}
\item \textbf{Space:} O($\frac{1}{\epsilon} \log(\epsilon n))$ tuples
\item \textbf{Update:} O($\frac{1}{\epsilon}$) amortized
\item \textbf{Query:} O($\frac{1}{\epsilon}$)
\item \textbf{Accuracy:} returns value within $\epsilon n$ of true rank
\end{itemize}

\subsection{Applications}

\begin{itemize}
\item \textbf{Monitoring dashboards:} P50, P90, P99 latency percentiles in real time
\item \textbf{Database query optimization:} estimate result distribution without scanning all rows
\item \textbf{Adaptive load balancing:} route requests based on recent response time distribution
\item \textbf{SLA monitoring:} are 99\% of requests served within 100ms?
\end{itemize}

\textbf{Alternatives:} t-digest (better accuracy for extreme quantiles), KLL sketch (better space-accuracy tradeoff). The GK sketch is the classic choice; t-digest is more common in production (used by Elasticsearch, Apache Flink).

\section{Reservoir Sampling}

Select k random elements from a stream of unknown length.

\begin{lstlisting}[language=C++]
std::vector<int> reservoir_sample(std::span<const int> stream, size_t k) {
    std::vector<int> reservoir(stream.begin(), stream.begin() + k);
    std::mt19937 rng(std::random_device{}());

    for (size_t i = k; i < stream.size(); ++i) {
        size_t j = std::uniform_int_distribution<size_t>(0, i)(rng);
        if (j < k) reservoir[j] = stream[i];
    }
    return reservoir;
}
\end{lstlisting}

Each element has probability k/n of being in the final sample.

\section{Sliding Window Model}

Most streaming algorithms assume you see the entire stream. In practice, you often need statistics over only the \textbf{last N elements} — a sliding window. The challenge: maintain a summary using O(polylog N) space while elements expire.

\subsection{Exponential Histogram}

The exponential histogram (Datar et al., 2002) maintains approximate counts over a sliding window of size W. It uses O($\frac{1}{\epsilon} \log W$) space to give $(1 \pm \epsilon)$-approximate counts.

The idea: divide the stream into buckets of exponentially increasing sizes. Each bucket covers a contiguous segment and stores its count. The smallest bucket covers 1 element, the next covers 2, then 4, 8, etc.

\begin{lstlisting}[language=C++]
template <typename T>
class sliding_window_counter {
public:
    explicit sliding_window_counter(size_t window_size, double epsilon = 0.01)
        : window_(window_size), epsilon_(epsilon) {}

    void add(const T& item, size_t timestamp) {
        // Add to bucket 0
        buckets_[0].push_back({item, 1});
        total_ += 1;

        // Merge buckets with same size if needed
        compact();

        // Remove expired elements
        while (!buckets_.empty() && buckets_.front().start_time +
               buckets_.front().count <= timestamp - window_) {
            total_ -= buckets_.front().count;
            buckets_.pop_front();
        }
    }

    size_t count() const { return total_; }

    // Approximate count of item in current window
    size_t approximate_count(const T& item) const {
        size_t result = 0;
        for (const auto& bucket : buckets_) {
            for (const auto& [val, cnt] : bucket.items) {
                if (val == item) result += cnt;
            }
        }
        return result;
    }

private:
    struct bucket {
        std::deque<std::pair<T, int>> items;
        size_t count = 0;
        size_t start_time = 0;
    };

    void compact() {
        // Merge consecutive buckets of same size
        for (size_t i = 0; i + 1 < buckets_.size(); ++i) {
            while (buckets_[i].count == buckets_[i+1].count &&
                   buckets_[i].count > 0) {
                // Merge bucket i+1 into bucket i
                for (auto& item : buckets_[i+1].items)
                    buckets_[i].items.push_back(item);
                buckets_[i].count += buckets_[i+1].count;
                buckets_.erase(buckets_.begin() + i + 1);
            }
        }
    }

    size_t window_;
    double epsilon_;
    std::deque<bucket> buckets_;
    size_t total_ = 0;
};
\end{lstlisting}

\subsection{Applications}

\begin{itemize}
\item \textbf{Network monitoring:} count packets/bytes in the last N seconds
\item \textbf{Rate limiting:} sliding window counters for API rate limits (Redis, NGINX)
\item \textbf{Log analysis:} error rate over the last hour, not the entire log
\item \textbf{Sensor networks:} recent average temperature, not lifetime average
\end{itemize}

\textbf{Complexity:} O($\frac{1}{\epsilon} \log W$) space, O(1) amortized per update.

\section{Consistent Hashing}

Distribute keys across a changing set of servers while minimizing remapping when servers join or leave. Each server is mapped to points on a hash ring, and each key is assigned to the nearest clockwise server.

\begin{lstlisting}[language=C++]
class consistent_hash_ring {
public:
    void add_server(std::string_view name, int virtual_nodes = 100) {
        for (int i = 0; i < virtual_nodes; ++i) {
            auto h = hash(std::string(name) + ":" + std::to_string(i));
            ring_[h] = name;
        }
    }

    void remove_server(std::string_view name) {
        std::erase_if(ring_, [&](const auto& p) {
            return p.second == name;
        });
    }

    std::string_view get_server(std::string_view key) const {
        if (ring_.empty()) return "";
        auto it = ring_.lower_bound(hash(key));
        if (it == ring_.end()) it = ring_.begin();
        return it->second;
    }

private:
    static uint64_t hash(std::string_view s) {
        return std::hash<std::string_view>{}(s);
    }

    std::map<uint64_t, std::string> ring_;
};
\end{lstlisting}

When a server is removed, only its keys migrate to the next server; every other key stays put. Virtual nodes improve load balance by spreading each server across many ring positions.

\section{Locality-Sensitive Hashing}

Approximate nearest-neighbor search in high dimensions by hashing inputs so that similar items collide with high probability.

The basic LSH\index{LSH} family for cosine similarity (simhash):
\begin{itemize}
\item Generate k random hyperplanes
\item For each item, compute the sign of each hyperplane dot product — the k-bit signature captures the angular neighborhood
\end{itemize}

\begin{lstlisting}[language=C++]
class simhash {
public:
    simhash(size_t bits, size_t dim)
        : bits_(bits), planes_(bits, std::vector<double>(dim)) {
        std::mt19937 rng(std::random_device{}());
        std::normal_distribution<double> dist(0.0, 1.0);
        for (auto& v : planes_)
            for (auto& x : v) x = dist(rng);
    }

    uint64_t signature(std::span<const double> vec) const {
        uint64_t sig = 0;
        for (size_t i = 0; i < bits_; ++i) {
            double dot = 0.0;
            for (size_t j = 0; j < vec.size(); ++j)
                dot += planes_[i][j] * vec[j];
            if (dot > 0) sig |= (1ULL << i);
        }
        return sig;
    }

    // Hamming distance on signatures approximates cosine distance
    static double cosine_similarity(uint64_t a, uint64_t b) {
        return 1.0 - static_cast<double>(std::popcount(a ^ b)) / 64.0;
    }

private:
    size_t bits_;
    std::vector<std::vector<double>> planes_;
};
\end{lstlisting}

\section{Merkle Tree (Hash Tree)}

A \textbf{Merkle tree} is a binary tree where each leaf stores the hash of a data block and each internal node stores the hash of its two children. It enables efficient verification that a specific block belongs to a large dataset without downloading the entire dataset.

\subsection{Structure}

\begin{lstlisting}[language=Java]
                  Hash(A||B)
                 /          \
            Hash(A)       Hash(B)
           /      \      /      \
     Hash(d0) Hash(d1) Hash(d2) Hash(d3)
       |        |        |        |
      d0       d1       d2       d3
\end{lstlisting}

\noindent The \textbf{root hash} uniquely represents the entire dataset. Changing any single data block changes the root hash, making tampering immediately detectable.

\subsection{Operations}

\begin{lstlisting}[language=C++]
#include <vector>
#include <string>
#include <cstdint>
#include <functional>
#include <stdexcept>

class merkle_tree {
public:
    using hash_fn = std::function<std::string(const std::string&)>;

    merkle_tree(std::vector<std::string> data, hash_fn h)
        : hash_(std::move(h)), data_(std::move(data)) {
        size_t n = data_.size();
        if (n == 0)
            throw std::invalid_argument("empty data");
        // Round up to next power of 2
        size_t leaves = 1;
        while (leaves < n) leaves *= 2;
        tree_.resize(2 * leaves, "");
        // Compute leaf hashes
        for (size_t i = 0; i < leaves; ++i) {
            if (i < n)
                tree_[leaves + i] = hash_(data_[i]);
            else
                tree_[leaves + i] = hash_("");  // padding
        }
        // Build internal nodes bottom-up
        for (size_t i = leaves - 1; i >= 1; --i)
            tree_[i] = hash_(tree_[2 * i] + tree_[2 * i + 1]);
        leaves_ = leaves;
    }

    std::string get_root_hash() const {
        return tree_[1];
    }

    // Merkle proof: list of sibling hashes from leaf to root
    std::vector<std::pair<std::string, bool>>
    get_proof(size_t index) const {
        std::vector<std::pair<std::string, bool>> proof;
        size_t pos = leaves_ + index;
        while (pos > 1) {
            bool is_left = (pos % 2 == 1);
            size_t sibling = is_left ? pos - 1 : pos + 1;
            proof.emplace_back(tree_[sibling], is_left);
            pos /= 2;
        }
        return proof;
    }

    // Verify that data_block is at index using the proof
    static bool verify_proof(
        hash_fn h,
        const std::string& data_block,
        size_t index,
        const std::string& expected_root,
        const std::vector<std::pair<std::string, bool>>& proof)
    {
        std::string current = h(data_block);
        for (const auto& [sibling_hash, is_left] : proof) {
            if (is_left)
                current = h(sibling_hash + current);
            else
                current = h(current + sibling_hash);
        }
        return current == expected_root;
    }

private:
    hash_fn hash_;
    std::vector<std::string> data_;
    std::vector<std::string> tree_;
    size_t leaves_;
};
\end{lstlisting}

\subsection{Complexity}

\begin{center}
\begin{tabular}{ll}
\toprule
Operation & Cost \\
\midrule
Build tree & $O(n)$ \\
Root hash & $O(1)$ \\
Generate proof & $O(\log n)$ \\
Verify proof & $O(\log n)$ \\
\bottomrule
\end{tabular}
\end{center}

\subsection{Applications}

\begin{itemize}
    \item \textbf{Blockchain:} Bitcoin and Ethereum use Merkle trees to summarize all transactions in a block. A light client can verify a single transaction with a $\log n$ proof instead of downloading every transaction.
    \item \textbf{Git:} Git commit objects contain a Merkle tree of the repository file tree, enabling efficient diff and integrity verification.
    \item \textbf{Distributed databases:} Cassandra and DynamoDB use Merkle trees to detect and reconcile inconsistencies between replicas during anti-entropy sync.
\end{itemize}

\section{Quotient Filter}

A \textbf{quotient filter}\index{quotient filter} is a space-efficient hash table that stores fingerprints using quotienting. Unlike a Bloom filter, it supports \textbf{counting and deletion}. A key $k$ is mapped to a fingerprint $f$ using a hash function, then split into quotient $q = f / 2^r$ and remainder $r = f \bmod 2^r$.

\subsection{Structure}

\begin{lstlisting}[language=Java]
  Quotient  |  Remainder slots
  ---------+-------------------------------
     0     |  [  _ ] [  _ ] [  _ ] [  _ ]
     1     |  [ b1 ] [  _ ] [  _ ] [  _ ]
     2     |  [ a0 ] [ a1 ] [  _ ] [  _ ]
     3     |  [ c0 ] [  _ ] [  _ ] [  _ ]
      .    |       .    .    .    .
\end{lstlisting}

\noindent Collisions with the same quotient form a \textbf{run} stored in consecutive slots using linear probing. A small auxiliary array tracks the start of each run.

\subsection{Operations}

\begin{lstlisting}[language=C++]
#include <vector>
#include <cstdint>
#include <functional>
#include <optional>

class quotient_filter {
public:
    quotient_filter(size_t capacity, size_t fingerprint_bits = 8)
        : r_(fingerprint_bits),
          size_(1ULL << (capacity > 0 ? compute_q_bits(capacity, fingerprint_bits) : 4)),
          mask_(size_ - 1),
          fingerprints_(size_, 0),
          occupied_(size_, 0),
          continuations_(size_, 0),
          counts_(size_, 0),
          size__(0)
    {
        q_bits_ = compute_q_bits(capacity, fingerprint_bits);
    }

    bool insert(uint64_t key) {
        uint64_t f = fingerprint(key);
        uint64_t q = f >> r_;
        uint64_t rem = f & ((1ULL << r_) - 1);
        if (counts_[q] == 0) {
            // Empty cluster -- place directly
            size_t pos = q;
            while (occupied_[pos]) pos = (pos + 1) & mask_;
            place(pos, q, rem);
            return true;
        }
        // Find the run for quotient q
        size_t run_start = find_run_start(q);
        // Check for duplicate
        size_t pos = run_start;
        for (size_t i = 0; i < counts_[q]; ++i) {
            if (fingerprints_[pos] == rem) return false;  // already present
            pos = (pos + 1) & mask_;
        }
        // Insert at end of run
        place(pos, q, rem);
        return true;
    }

    bool lookup(uint64_t key) const {
        uint64_t f = fingerprint(key);
        uint64_t q = f >> r_;
        uint64_t rem = f & ((1ULL << r_) - 1);
        if (counts_[q] == 0) return false;
        size_t pos = find_run_start(q);
        for (size_t i = 0; i < counts_[q]; ++i) {
            if (fingerprints_[pos] == rem) return true;
            pos = (pos + 1) & mask_;
        }
        return false;
    }

    bool remove(uint64_t key) {
        uint64_t f = fingerprint(key);
        uint64_t q = f >> r_;
        uint64_t rem = f & ((1ULL << r_) - 1);
        if (counts_[q] == 0) return false;
        size_t pos = find_run_start(q);
        for (size_t i = 0; i < counts_[q]; ++i) {
            if (fingerprints_[pos] == rem) {
                // Shift run left and compact
                compact(pos, q);
                return true;
            }
            pos = (pos + 1) & mask_;
        }
        return false;
    }

    size_t size() const { return size__; }

private:
    size_t r_, q_bits_, size_, mask_;
    size_t size__;
    std::vector<uint64_t> fingerprints_;
    std::vector<bool> occupied_;
    std::vector<bool> continuations_;
    std::vector<size_t> counts_;

    static size_t compute_q_bits(size_t cap, size_t r) {
        size_t total_slots = cap * 2;  // 50% load factor
        size_t q = 0;
        while ((1ULL << q) < total_slots) ++q;
        return q;
    }

    uint64_t fingerprint(uint64_t key) const {
        // SplitMix64-style hash
        key ^= key >> 30;
        key *= 0xbf58476d1ce4e5b9ULL;
        key ^= key >> 27;
        key *= 0x94d049bb133111ebULL;
        key ^= key >> 31;
        return key & ((1ULL << (q_bits_ + r_)) - 1);
    }

    size_t find_run_start(uint64_t q) const {
        // Count occupied slots before q
        size_t count = 0;
        for (size_t i = 0; i < q; ++i)
            if (occupied_[i]) ++count;
        // Scan from q to find the (count+1)-th occupied cluster
        size_t pos = q;
        size_t seen = 0;
        while (seen < count) {
            pos = (pos + 1) & mask_;
            if (occupied_[pos]) ++seen;
        }
        return (pos + 1) & mask_;
    }

    void place(size_t pos, uint64_t q, uint64_t rem) {
        fingerprints_[pos] = rem;
        occupied_[q] = true;
        ++counts_[q];
        ++size__;
    }

    void compact(size_t pos, uint64_t q) {
        // Shift elements left to fill the gap
        size_t next = (pos + 1) & mask_;
        while (counts_[q] > 1 || (next != pos && occupied_[next])) {
            fingerprints_[pos] = fingerprints_[next];
            pos = next;
            next = (next + 1) & mask_;
        }
        // Clear the last slot
        occupied_[q] = false;
        if (counts_[q] > 0) --counts_[q];
        --size__;
    }
};
\end{lstlisting}

\subsection{Complexity}

\begin{center}
\begin{tabular}{ll}
\toprule
Operation & Cost \\
\midrule
Insert & $O(1)$ amortized \\
Lookup & $O(1)$ amortized \\
Delete & $O(1)$ amortized \\
Space & $O(n)$ fingerprints \\
\bottomrule
\end{tabular}
\end{center}

\subsection{Comparison with Bloom Filter}

\begin{center}
\begin{tabular}{lcc}
\toprule
Property & Bloom Filter & Quotient Filter \\
\midrule
Supports deletion & No & Yes \\
Space efficiency & Moderate & Higher (smaller fingerprint) \\
Cache performance & Poor & Good (sequential probing) \\
False positives & Yes & Yes (via hash collision) \\
\bottomrule
\end{tabular}
\end{center}

\noindent The quotient filter's sequential memory layout gives it better cache locality than a Bloom filter of equivalent size, making it preferred when deletion support is needed.

\section{MinHash}

\textbf{MinHash}\index{MinHash} estimates the \textbf{Jaccard similarity} between two sets without comparing all elements. Given sets $A$ and $B$, the Jaccard similarity is $J(A,B) = |A \cap B| / |A \cup B|$. MinHash uses $k$ independent hash functions; for each function, the minimum hash value over a set is computed, producing a \textbf{signature} of $k$ values.

\subsection{Structure}

\begin{lstlisting}[language=Java]
Set A = {a, c, d, f, g}    Set B = {b, c, e, f, g}

  h1:  min(h1(A)) = h1(f)     min(h1(B)) = h1(b)
  h2:  min(h2(A)) = h2(c)     min(h2(B)) = h2(c)
  h3:  min(h3(A)) = h3(d)     min(h3(B)) = h3(e)

  Sig(A) = [h1(f), h2(c), h3(d)]
  Sig(B) = [h1(b), h2(c), h3(e)]
  Matches: 1 out of 3 => J estimate = 0.33
\end{lstlisting}

\noindent For each hash function $h_i$, $P[\min(h_i(A)) = \min(h_i(B))] = J(A,B)$. The fraction of matching entries in two signatures is an unbiased estimator of the Jaccard similarity.

\subsection{Operations}

\begin{lstlisting}[language=C++]
#include <vector>
#include <unordered_set>
#include <cstdint>
#include <random>
#include <algorithm>
#include <numeric>

class minhash {
public:
    minhash(size_t num_hashes, size_t seed = 42)
        : k_(num_hashes), hashes_(num_hashes)
    {
        // Generate k independent hash function parameters
        std::mt19937_64 rng(seed);
        for (size_t i = 0; i < k_; ++i) {
            uint64_t a = rng(), b = rng();
            hashes_[i] = [a, b](uint64_t x) -> uint64_t {
                // 64-bit multiply-shift hash
                return a * x + b;
            };
        }
    }

    // Compute MinHash signature for a set of elements
    std::vector<uint64_t> compute_signature(
        const std::unordered_set<std::string>& set) const
    {
        std::vector<uint64_t> sig(k_, UINT64_MAX);
        for (const auto& elem : set) {
            uint64_t h = std::hash<std::string>{}(elem);
            for (size_t i = 0; i < k_; ++i) {
                uint64_t val = hashes_[i](h);
                sig[i] = std::min(sig[i], val);
            }
        }
        return sig;
    }

    // Estimate Jaccard similarity from two signatures
    static double estimate_jaccard(
        const std::vector<uint64_t>& sig_a,
        const std::vector<uint64_t>& sig_b)
    {
        size_t matches = 0;
        size_t k = sig_a.size();
        for (size_t i = 0; i < k; ++i)
            if (sig_a[i] == sig_b[i]) ++matches;
        return static_cast<double>(matches) / k;
    }

    size_t num_hashes() const { return k_; }

private:
    size_t k_;
    std::vector<std::function<uint64_t(uint64_t)>> hashes_;
};
\end{lstlisting}

\subsection{Complexity}

\begin{center}
\begin{tabular}{ll}
\toprule
Operation & Cost \\
\midrule
Compute signature & $O(k \cdot n)$ where $n = |S|$ \\
Estimate Jaccard & $O(k)$ \\
Signature size & $O(k)$ \\
\bottomrule
\end{tabular}
\end{center}

\noindent Increasing $k$ reduces the variance of the estimate. With $k = 200$ the standard error is approximately $1/\sqrt{k} \approx 0.07$, sufficient for most similarity tasks.

\subsection{Applications}

\begin{itemize}
    \item \textbf{Near-duplicate detection:} Web crawlers use MinHash to detect pages with high Jaccard similarity, filtering redundant content at scale.
    \item \textbf{Recommendation systems:} Users with similar item sets (high Jaccard) receive similar recommendations; MinHash makes this comparison sublinear.
    \item \textbf{Locality-sensitive hashing\index{locality-sensitive hashing} (LSH):} MinHash signatures are banded and hashed to find \emph{candidate pairs} in $O(n)$ expected time, enabling subquadratic all-pairs similarity search.
\end{itemize}

\section{HyperLogLog}

HyperLogLog\index{HyperLogLog} estimates the number of distinct elements in a stream using $O(m)$ registers (typically $m = 2^p$ for $10 \le p \le 16$). The key insight is that a uniformly distributed hash produces leading zeros with probability $2^{-j}$; the maximum number of leading zeros across $m$ buckets gives a raw estimate, which is then corrected using the harmonic mean.

\subsection{Core Idea}

Given a 64-bit hash function, for each element:
\begin{enumerate}
    \item Hash the element: $h = \text{hash}(x)$.
    \item Use the first $p$ bits to select one of $m = 2^p$ buckets.
    \item In the remaining $64 - p$ bits, count the leading zeros $\rho(w)$ (plus 1).
    \item Store $\max(\rho(w))$ in the selected bucket.
\end{enumerate}

The raw cardinality estimate is:
\[
E = \alpha_m \cdot m^2 \cdot \left( \sum_{j=1}^{m} 2^{-M[j]} \right)^{-1}
\]
where $M[j]$ is the register value for bucket $j$ and $\alpha_m$ is a bias-correction constant. For $m \ge 128$, $\alpha_m \approx 0.7213 / (1 + 1.079 / m)$.

A linear counting correction is applied when $E \le 2.5 \cdot m$: count the number of zero registers $V$ and use $E' = m \cdot \ln(m / V)$ instead.

\subsection{Standard Error}

The standard error of HyperLogLog is approximately $1.04 / \sqrt{m}$. With $p = 14$ ($m = 16384$ registers, 20 KB), the standard error is $\approx 0.81\%$, giving approximately 2.5\% accuracy at the 95\% confidence level.

\begin{lstlisting}[language=C++, caption={HyperLogLog cardinality estimator}]
#include <array>
#include <cstdint>
#include <cmath>
#include <algorithm>
#include <bit>

class HyperLogLog {
public:
    explicit HyperLogLog(uint8_t p)
        : p_(p), m_(1u << p), alpha_(0.7213 / (1.0 + 1.079 / m_))
    {
        M_.fill(0);
    }

    void add(uint64_t hash) {
        uint32_t idx = hash >> (64 - p_);
        uint64_t w = hash << p_ | (1ULL << (p_ - 1));
        uint8_t rho = static_cast<uint8_t>(
            std::countl_zero(w) + 1);
        M_[idx] = std::max(M_[idx], rho);
    }

    double estimate() const {
        double sum = 0.0;
        int zeros = 0;
        for (int j = 0; j < m_; ++j) {
            sum += std::pow(2.0, -M_[j]);
            if (M_[j] == 0) ++zeros;
        }
        double E = alpha_ * m_ * m_ / sum;
        if (E <= 2.5 * m_ && zeros > 0) {
            E = m_ * std::log(static_cast<double>(m_) / zeros);
        }
        if (E > (1ULL << 32) / 30.0) {
            E = -std::pow(2.0, 32) *
                std::log1p(-E / std::pow(2.0, 32));
        }
        return E;
    }

    void merge(const HyperLogLog& other) {
        for (int j = 0; j < m_; ++j)
            M_[j] = std::max(M_[j], other.M_[j]);
    }

    size_t memory_bytes() const {
        return m_ + 2 * sizeof(uint8_t) + sizeof(int);
    }

private:
    uint8_t p_;
    int m_;
    double alpha_;
    std::array<uint8_t, 1 << 16> M_;
};
\end{lstlisting}

\subsection{Complexity}

\begin{center}
\begin{tabular}{ll}
\toprule
Operation & Cost \\
\midrule
Add element & $O(1)$ \\
Estimate cardinality & $O(m)$ \\
Merge two structures & $O(m)$ \\
Memory & $O(m)$ bytes (typically 2--64 KB) \\
\bottomrule
\end{tabular}
\end{center}

\subsection{Applications}
\begin{itemize}
    \item \textbf{Network monitoring:} Estimate unique IP addresses in a time window without storing them all.
    \item \textbf{Database query planning:} Cardinality estimates for JOIN and GROUP BY operations drive cost-based optimizers.
    \item \textbf{Analytics:} Count unique users, unique sessions, or unique search queries in large-scale event streams.
\end{itemize}

\section{Reservoir Sampling}

Reservoir sampling selects a uniform random sample of $k$ elements from a stream of $n$ elements where $n$ is unknown or too large to hold in memory. Algorithm R (Vitter, 1985) is the simplest approach.

\subsection{Algorithm R}

Maintain a reservoir $R$ of size $k$. For the $i$-th element ($i \ge 1$):
\begin{enumerate}
    \item If $i \le k$, insert the element into $R$.
    \item If $i > k$, generate a random integer $j$ uniformly from $\{1, 2, \ldots, i\}$. If $j \le k$, replace $R[j-1]$ with the new element.
\end{enumerate}

The probability that any element ends up in the sample is exactly $k / n$, proven by induction on $i$.

\subsection{Correctness Argument}

At step $i > k$:
\[
\Pr[\text{element } i \text{ is in } R] = \frac{k}{i}
\]
For element $j \le i$ to survive step $i$:
\[
\Pr[\text{element } j \text{ survives step } i] = 1 - \frac{1}{i} = \frac{i-1}{i}
\]
After $n$ steps, the survival probability is:
\[
\prod_{i=j+1}^{n} \frac{i-1}{i} = \frac{j}{n}
\]
Thus each element has probability $k / n$ of being selected.

\begin{lstlisting}[language=C++, caption={Reservoir sampling (Algorithm R)}]
#include <vector>
#include <random>
#include <iterator>

template <typename Iter>
std::vector<typename std::iterator_traits<Iter>::value_type>
reservoir_sample(Iter begin, Iter end, size_t k,
                 std::mt19937& rng)
{
    std::vector<typename std::iterator_traits<Iter>::value_type>
        R;
    R.reserve(k);

    size_t i = 0;
    for (auto it = begin; it != end; ++it, ++i) {
        if (i < k) {
            R.push_back(*it);
        } else {
            std::uniform_int_distribution<size_t>
                dist(0, i);
            size_t j = dist(rng);
            if (j < k)
                R[j] = *it;
        }
    }
    return R;
}

template <typename T>
class StreamSampler {
public:
    explicit StreamSampler(size_t k, uint64_t seed = 42)
        : k_(k), rng_(seed), n_(0) {
        R_.reserve(k);
    }

    void add(const T& value) {
        if (n_ < k_) {
            R_.push_back(value);
        } else {
            std::uniform_int_distribution<size_t>
                dist(0, n_);
            size_t j = dist(rng_);
            if (j < k_)
                R_[j] = value;
        }
        ++n_;
    }

    const std::vector<T>& sample() const { return R_; }
    size_t total_count() const { return n_; }

private:
    size_t k_;
    std::mt19937 rng_;
    uint64_t n_;
    std::vector<T> R_;
};
\end{lstlisting}

\subsection{Complexity}

\begin{center}
\begin{tabular}{ll}
\toprule
Operation & Cost \\
\midrule
Add element & $O(1)$ \\
Memory & $O(k)$ \\
Space & $O(k)$ words \\
\bottomrule
\end{tabular}
\end{center}

\subsection{Applications}
\begin{itemize}
    \item \textbf{Data analytics:} Sampling from massive datasets for A/B testing or exploratory analysis.
    \item \textbf{Stream processing:} Selecting representative queries from a log stream for debugging.
    \item \textbf{Machine learning:} Creating balanced training sets from imbalanced streams.
\end{itemize}

\section{Count-Distinct (Linear Counting)}

Linear counting is the simplest distinct-element counting technique: allocate a bitmap of $m$ bits, hash each element to a position, and set that bit. The number of zero bits $V$ gives an estimate of the distinct count $n$.

\subsection{Estimate Formula}

The fraction of zero bits follows:
\[
\frac{V}{m} \approx e^{-n/m}
\]
Rearranging gives the linear counting estimate:
\[
\hat{n} = -m \cdot \ln\left(\frac{V}{m}\right)
\]
This estimate is accurate when the bitmap is not too full (less than $\sim 70\%$ set). Once $V / m$ drops below a threshold, switch to another method or double $m$.

\begin{lstlisting}[language=C++, caption={Linear counting for distinct elements}]
#include <vector>
#include <functional>
#include <cmath>

class LinearCounter {
public:
    explicit LinearCounter(size_t m)
        : m_(m), bits_(m, false), count_(0) {}

    void add(uint64_t hash) {
        size_t idx = hash % m_;
        if (!bits_[idx]) {
            bits_[idx] = true;
            ++count_;
        }
    }

    double estimate() const {
        if (count_ == m_) return static_cast<double>(m_);
        double V = static_cast<double>(m_ - count_);
        return -m_ * std::log(V / m_);
    }

    double load_factor() const {
        return static_cast<double>(count_) / m_;
    }

    size_t memory_bytes() const {
        return m_ / 8 + 2 * sizeof(size_t);
    }

private:
    size_t m_;
    std::vector<bool> bits_;
    size_t count_;
};
\end{lstlisting}

\subsection{When to Use}

\begin{itemize}
    \item Use when cardinality is expected to be small relative to $m$ (bitmap not too full).
    \item Provides a simple, fast baseline; HyperLogLog is preferred for large cardinalities.
    \item Linear counting uses roughly $n / \ln 2 \approx 1.44n$ bits when $V \approx m / 2$.
\end{itemize}

\section{Exercises}

\begin{enumerate}
\item Implement HyperLogLog with $p = 10, 12, 14, 16$ and estimate the number of distinct values in a stream of $10^7$ uniform random integers. Plot standard error vs.\ $p$ and compare with the theoretical value $1.04 / \sqrt{2^p}$.
\end{enumerate}

\begin{enumerate}
\item Implement a streaming system that switches from linear counting to HyperLogLog when the bitmap load factor exceeds 70\%. Compare accuracy at the switch point with using HyperLogLog alone from the start.
\end{enumerate}

\begin{enumerate}
\item Use reservoir sampling with $k = 50$ to sample from a stream of $10^6$ integers where each integer $x$ appears with probability proportional to $1 / \log(x + 2)$. Verify empirically that the frequency of each sampled element matches its true frequency within expected error.
\end{enumerate}

\begin{enumerate}
\item Merge two HyperLogLog structures with different seeds on disjoint halves of a dataset. Does the merged estimate match the true distinct count? How does accuracy degrade when the two halves overlap by 50\%?
\end{enumerate}

\begin{enumerate}
\item Implement a probabilistic counter using $b$ bitmaps of size $m = 8192$ with different hash seeds. At query time, use the bitmap with the lowest load factor to estimate cardinality. Compare accuracy and space with a single large bitmap of size $b \cdot m$.
\end{enumerate}

\subsection{Top-K Approximation}

Finding the top-$k$ most frequent elements in a data stream is a common problem
in network traffic analysis, trending topics, and log monitoring.
A min-heap of size $k$ paired with a hash map provides an online algorithm that
approximates the top-$k$ elements without storing the entire stream.

\begin{lstlisting}[language=C++]
#include <iostream>
#include <queue>
#include <unordered_map>
#include <vector>
#include <string>

using Entry = std::pair<int, std::string>;   // (count, element)

// Min-heap of size k: keep the k elements with the largest counts.
class TopK {
public:
    TopK(std::size_t k) : k_(k) {}

    void insert(const std::string& item) {
        int& cnt = freq_[item];
        cnt++;
        if (in_heap_.count(item)) {
            // Already in heap; will be rebuilt on next topK().
            stale_ = true;
            return;
        }
        if (heap_.size() < k_) {
            heap_.emplace(cnt, item);
            in_heap_.insert(item);
        } else if (cnt > heap_.top().first) {
            auto [old_cnt, old_item] = heap_.top();
            heap_.pop();
            in_heap_.erase(old_item);
            heap_.emplace(cnt, item);
            in_heap_.insert(item);
        }
    }

    std::vector<std::string> topK() {
        if (stale_) rebuild();
        std::vector<std::string> result;
        while (!heap_.empty()) {
            result.push_back(heap_.top().second);
            heap_.pop();
        }
        std::reverse(result.begin(), result.end());
        in_heap_.clear();
        stale_ = false;
        return result;
    }

private:
    void rebuild() {
        // Rebuild heap from current frequency map.
        std::priority_queue<Entry, std::vector<Entry>, std::greater<>> fresh;
        for (auto& [item, cnt] : freq_) {
            if (fresh.size() < k_) {
                fresh.emplace(cnt, item);
            } else if (cnt > fresh.top().first) {
                fresh.pop();
                fresh.emplace(cnt, item);
            }
        }
        heap_ = std::move(fresh);
    }

    std::size_t k_;
    std::unordered_map<std::string, int> freq_;
    std::priority_queue<Entry, std::vector<Entry>, std::greater<>> heap_;
    std::unordered_set<std::string> in_heap_;
    bool stale_ = false;
};
\end{lstlisting}

The hash map accumulates exact counts; the min-heap evicts the smallest count
when a new candidate exceeds the heap's minimum.  When a heap-resident element
receives an update, the algorithm marks the heap as stale and rebuilds it on
demand, keeping the structure correct without per-update re-heapification.
The space cost is $O(k + u)$, where $u$ is the number of distinct items seen and
$k$ is the desired output size.

\section{Common Bugs and Pitfalls}

\begin{center}
\begin{tabular}{ccc}
\toprule
Pitfall & Consequence & Solution \\
Bloom filter false positive rate exceeds requirements & Too many false positives, wasted queries & Calculate optimal k and m from n and desired P; verify with test data \\
Using Bloom filter when deletions are needed & Standard Bloom filter doesn't support deletion & Use counting Bloom filter or switch to a Cuckoo filter \\
Cuckoo filter insertion failure (too full) & Insert returns without adding element & Rebuild with 2$\times $ buckets; or use counting Bloom filter \\
Cuckoo filter max kicks exceeded during eviction & Infinite loop or degraded performance & Limit evictions (500), trigger rebuild when exceeded \\
Forgetting Consistent Hashing's virtual nodes & Skewed distribution for non-uniform servers & Add 100–200 virtual nodes per physical server \\
Hash collision in MinHash / Simhash & Wrong similarity estimate & Use multiple hash functions (k = 100+ for MinHash) \\
HyperLogLog with too few registers (m too small) & High estimation error & Use m = 2$^{12}$ (4096 registers) for \textasciitilde{}1.6\% error \\
Using reservoir sampling without shuffling first n elements & First k elements overrepresented, bias & Implement Fisher-Yates for the first k positions \\
LSH hash table with too few bands & Poor recall (many near neighbors missed) & Tune band count b and row count r for desired (recall, precision) \\

\bottomrule
\end{tabular}
\end{center}
\section{Summary}

\begin{center}
\begin{tabular}{ccccc}
\toprule
Structure & Problem & Space & Error & Supports Deletion \\
Bloom filter & Membership & O(m) bits & False positive rate P & With counting variant \\
Cuckoo filter & Membership & O(m) bits & False positive rate P & Yes \\
Count-Min Sketch & Frequency & O(dw) counters & Additive $\varepsilon $$\cdot $n & No \\
Heavy Hitters (Misra-Gries) & Top-k items & O(k) counters & Exact for f $>$ n/k & No \\
HyperLogLog & Cardinality & O(m log log n) & 1.04/$\sqrt {m}$ relative & No \\
Quantile Estimation (GK) & Quantiles & O($\frac{1}{\varepsilon}$ log $\varepsilon $n) & $\varepsilon $n rank error & No \\
Reservoir Sampling & Random sample & O(k) & Uniform -- exact & No \\
Sliding Window & Window count & O($\frac{1}{\varepsilon}$ log W) & $(1 \pm \varepsilon)$ & Yes \\
Consistent Hashing & Key distribution & O(nV) mapping & Near-perfect balance & Yes \\
LSH / Simhash & Approx NN & O(k) per item & Depends on k & -- \\

\bottomrule
\end{tabular}
\end{center}
These trade-offs are essential for streaming data, big data analytics, distributed systems, and ML pipelines.

\section{Exercises}

\subsection{Drill}

\begin{enumerate}
\item For a Bloom filter with m = 100 bits, k = 3 hash functions, and n = 10 elements inserted, what is the expected false positive rate? Compute using the formula (1 - e\textasciicircum{}{-kn/m})\textasciicircum{}k.
\end{enumerate}

\begin{enumerate}
\item How many bits are needed for a Bloom filter to achieve a false positive rate of 1\% with 10$^{6}$ elements? What if we need 0.1\%?
\end{enumerate}

\begin{enumerate}
\item For Count-Min Sketch, if we want to guarantee that all frequency estimates have error at most 0.1\% with probability 99\%, how many rows and columns do we need?
\end{enumerate}

\begin{enumerate}
\item For HyperLogLog with 14-bit precision (m = 16384 registers), what is the standard error? If we halve the precision (use 13 bits), how does the error change?
\end{enumerate}

\begin{enumerate}
\item Explain why Bloom filters have false positives but no false negatives. Give a concrete example of a false positive.
\end{enumerate}

\begin{enumerate}
\item A cuckoo filter with bucket size 4 and 8 bits per fingerprint has a 1\% false positive rate. How many buckets are needed for 10$^{6}$ elements? How does this compare with a Bloom filter using the same false positive rate?
\end{enumerate}

\subsection{Application}

\begin{enumerate}
\item Implement a \textbf{counting Bloom filter} that supports deletion by replacing each bit with a small counter (typically 4 bits). Compare space with the standard Bloom filter.
\end{enumerate}

\begin{enumerate}
\item Implement \textbf{Bloom filter union and intersection} — given two Bloom filters with the same parameters, compute the union (OR) and intersection (AND). Verify that the union correctly represents the union of the two sets.
\end{enumerate}

\begin{enumerate}
\item Implement a \textbf{Bloom filter vs Cuckoo filter benchmark}: insert 10$^{6}$ elements, then measure space, insertion throughput, lookup throughput, and deletion throughput for each. At what false positive rate does the Cuckoo filter become strictly better?
\end{enumerate}

\begin{enumerate}
\item Compare Count-Min Sketch vs \texttt{std::unordered\_map} for frequency estimation\index{frequency estimation} on a stream of 10$^{6}$ items with 10$^{5}$ distinct keys. Measure the error in the sketch's estimates and the memory used.
\end{enumerate}

\begin{enumerate}
\item Use HyperLogLog to estimate the number of distinct IP addresses in a network log with 10$^{7}$ entries. Compare accuracy at precision values p = 10, 12, 14, 16.
\end{enumerate}

\begin{enumerate}
\item Implement \textbf{reservoir sampling} for k = 100 from a stream of unknown length. Verify that each element has probability k/n of being in the sample through simulation.
\end{enumerate}

\begin{enumerate}
\item Implement a \textbf{heavy hitters} algorithm (lossy counting or space-saving) that finds all elements with frequency > $\phi $$\cdot $n in a stream. Compare with Count-Min Sketch.
\end{enumerate}

\begin{enumerate}
\item Implement a \textbf{Bloom filter} for a web crawler to avoid revisiting URLs. How does the false positive rate affect the number of missed pages?
\end{enumerate}

\begin{enumerate}
\item Implement \textbf{consistent hashing} for a distributed cache with 5 servers and 100 keys. Show the distribution of keys when a server joins and when one leaves. Use 50 virtual nodes per server.
\end{enumerate}

\begin{enumerate}
\item Use \textbf{simhash} to find near-duplicate documents in a collection of 100 text snippets. Compute the cosine similarity from signatures and compare with exact cosine similarity of TF-IDF vectors.
\end{enumerate}

\subsection{Research}

\begin{enumerate}
\item \textbf{(MinHash)} Read about MinHash for estimating set similarity (Jaccard coefficient). Implement MinHash with k hash functions and compare with exact Jaccard computation for random sets.
\end{enumerate}

\begin{enumerate}
\item \textbf{(Cuckoo filter scaling)} Our cuckoo filter uses a fixed bucket size of 4 and a max of 500 kicks. Experiment with different bucket sizes (2, 4, 8) and max kick limits. How do these parameters affect the fill factor before a rebuild is needed? At what load factor does the filter become unstable?
\end{enumerate}

\begin{enumerate}
\item \textbf{(HyperLogLog++)} Read about Google's HyperLogLog++ improvements (Heule et al., 2013) including bias correction and sparse representation. Implement the sparse representation for small cardinalities.
\end{enumerate}

\begin{enumerate}
\item \textbf{(Count-Min Sketch variants)} Research the Conservative Update variant of Count-Min Sketch. Implement it and measure the improvement in accuracy.
\end{enumerate}

\begin{enumerate}
\item \textbf{(DDSketch)} Research the DDSketch algorithm for streaming quantiles. Implement it and compare accuracy with the Greenwald-Khanna algorithm on skewed distributions.
\end{enumerate}

\begin{enumerate}
\item \textbf{(Rendezvous hashing)} Implement HRW (highest random weight) hashing and compare load distribution with consistent hashing. Which achieves better balance with fewer virtual nodes?
\end{enumerate}

\section{References and Further Reading}

\begin{itemize}
\item Burton H. Bloom, "Space/Time Trade-offs in Hash Coding with Allowable Errors" — CACM, 1970.
\item Graham Cormode and S. Muthukrishnan, "An Improved Data Stream Summary: The Count-Min Sketch and Its Applications" — Journal of Algorithms, 2005.
\item Philippe Flajolet, Eric Fusy, Olivier Gandouet, and Frederic Meunier, "HyperLogLog: the Analysis of a Near-Optimal Cardinality Estimation Algorithm" — AOFA, 2007.
\item Jeffrey S. Vitter, "Random Sampling with a Reservoir" — ACM TOMS, 1985.
\item Andrei Z. Broder, "On the Resemblance and Containment of Documents" — Compression and Complexity of Sequences, 1997.
\item Bin Fan, Dave G. Andersen, Michael Kaminsky, and Michael D. Mitzenmacher, "Cuckoo Filter: Practically Better Than Bloom" — ACM CoNEXT, 2014.
\item Stefan Heule, Marc Nunkesser, and Alexander Hall, "HyperLogLog in Practice: Algorithmic Engineering of a State of the Art Cardinality Estimation Algorithm" — EDBT, 2013.
\item Jayadev Misra and David Gries, "Finding Repeated Elements" — Science of Computer Programming, 1982.
\item Michael Greenwald and Sanjeev Khanna, "Space-Efficient Online Computation of Quantile Summaries" — SIGMOD, 2001 (GK sketch).
\item Gurmeet Singh Manku and Rajeev Motwani, "Approximate Frequency Counts over Data Streams" — VLDB, 2002 (Misra-Gries and related algorithms).
\item Datar, Gionis, Indyk, and Motwani, "Maintaining Stream Statistics over Sliding Windows" — SIAM Review, 2002 (exponential histogram).
\item David Karger et al., "Consistent Hashing and Random Trees: Distributed Caching Protocols for Relieving Hot Spots on the World Wide Web" — STOC, 1997.
\item Piotr Indyk and Rajeev Motwani, "Approximate Nearest Neighbors: Towards Removing the Curse of Dimensionality" — STOC, 1998.
\item Moses S. Charikar, "Similarity Estimation Techniques from Rounding Algorithms" — STOC, 2002.
\item Ted Dunning, "The T-Digest: Efficient Estimates of Distributions" — 2013.
\item Graham Cormode, S. Muthukrishnan, Korn, and Fang, "Towards a Science of Data Streaming" — (survey of streaming algorithms).
\end{itemize}

